\chapter{Donn\'ees de sant\'e avec Pandas}
\label{ch:pandas}

\begin{quote}
\emph{<<\,80\,\% de l'analyse de donn\'ees consiste \`a mettre les donn\'ees dans une forme o\`u l'on peut leur poser des questions.\,>>}
\end{quote}

% ============================================================
\section{Qu'est-ce qu'un DataFrame ?}

Un DataFrame est un tableau --- les lignes sont des observations (patients, pays, points temporels), les colonnes sont des variables (\'age, diagnostic, valeurs de laboratoire). Si vous avez utilis\'e Excel, vous comprenez d\'ej\`a le concept. Pandas vous donne la puissance d'Excel avec la reproductibilit\'e de Python.

\begin{minted}{python}
import pandas as pd

# Creer un petit DataFrame clinique
data = {
    "patient_id": ["P001", "P002", "P003", "P004", "P005"],
    "age": [45, 67, 34, 52, 71],
    "sex": ["M", "F", "F", "M", "F"],
    "systolic_bp": [128, 158, 112, 145, 162],
    "hba1c": [5.4, 7.8, 5.1, 6.3, 8.2],
    "diagnosis": ["None", "T2DM", "None", "Pre-DM", "T2DM"]
}
df = pd.DataFrame(data)
df
\end{minted}

% ============================================================
\section{Charger des donn\'ees de sant\'e r\'eelles}

\subsection{Fichiers CSV}

La plupart des jeux de donn\'ees de sant\'e sont fournis sous forme de fichiers CSV (valeurs s\'epar\'ees par des virgules) :

\begin{minted}{python}
# Donnees d'esperance de vie de l'OMS (Gapminder)
url = "https://raw.githubusercontent.com/datasets/gapminder/main/data/gapminder.csv"
gapminder = pd.read_csv(url)
print(f"Dimensions : {gapminder.shape}")  # (lignes, colonnes)
gapminder.head()
\end{minted}

\subsection{Fichiers Excel}

\begin{minted}{python}
df = pd.read_excel("patient_records.xlsx", sheet_name="Lab Results")
\end{minted}

\subsection{Depuis l'Observatoire mondial de la sant\'e de l'OMS}

\begin{minted}{python}
# Ratio de mortalite maternelle par pays
url = "https://apps.who.int/gho/athena/api/GHO/MDG_0000000026?format=csv"
mmr = pd.read_csv(url)
mmr.head()
\end{minted}

\begin{clinicalnote}
L'API GHO de l'OMS fournit plus de 1\,000 indicateurs de sant\'e pour 194 \'Etats membres. Vous pouvez parcourir les indicateurs sur \url{https://www.who.int/data/gho/data/indicators}. Chaque indicateur poss\`ede un code (par ex.\ \texttt{MDG\_0000000026} pour le ratio de mortalit\'e maternelle) que vous int\'egrez dans l'URL de l'API.
\end{clinicalnote}

% ============================================================
\section{Explorer vos donn\'ees}

La premi\`ere chose \`a faire avec tout jeu de donn\'ees de sant\'e --- avant toute analyse --- est de \emph{le regarder}.

\begin{minted}{python}
# Exploration de base
print(gapminder.shape)          # dimensions
print(gapminder.columns.tolist())  # noms des colonnes
print(gapminder.dtypes)         # types de donnees
gapminder.describe()            # statistiques descriptives
gapminder.info()                # comptage des non-nuls et memoire
\end{minted}

\subsection{S\'electionner des colonnes}

\begin{minted}{python}
# Une seule colonne (retourne une Series)
ages = gapminder["lifeExp"]

# Plusieurs colonnes (retourne un DataFrame)
subset = gapminder[["country", "year", "lifeExp"]]
\end{minted}

\subsection{Filtrer des lignes}

\begin{minted}{python}
# Patients avec HbA1c >= 6.5 (seuil diabetique)
diabetic = df[df["hba1c"] >= 6.5]

# Pays africains dans Gapminder
africa = gapminder[gapminder["continent"] == "Africa"]

# Conditions multiples : patientes de plus de 60 ans avec PA elevee
high_risk = df[(df["sex"] == "F") & (df["age"] > 60) & (df["systolic_bp"] >= 140)]
\end{minted}

\begin{warningbox}
Utilisez \texttt{\&} (et), \texttt{|} (ou), \texttt{\textasciitilde} (non) avec des parenth\`eses autour de chaque condition. Les mots-cl\'es Python \texttt{and}/\texttt{or} ne fonctionnent pas avec les DataFrames.
\end{warningbox}

% ============================================================
\section{Tri et classement}

\begin{minted}{python}
# Pays avec l'esperance de vie la plus elevee en 2007
recent = gapminder[gapminder["year"] == 2007]
top10 = recent.sort_values("lifeExp", ascending=False).head(10)
print(top10[["country", "lifeExp"]])
\end{minted}

% ============================================================
\section{Regroupement et agr\'egation}

Le regroupement est la fa\c{c}on de r\'epondre \`a des questions comme <<\,quelle est l'esp\'erance de vie moyenne par continent ?\,>> ou <<\,combien de patients dans chaque cat\'egorie de diagnostic ?\,>>

\begin{minted}{python}
# Esperance de vie moyenne par continent (2007)
recent.groupby("continent")["lifeExp"].mean().sort_values(ascending=False)
\end{minted}

\begin{minted}{python}
# Agregations multiples
recent.groupby("continent")["lifeExp"].agg(["mean", "median", "std", "count"])
\end{minted}

\begin{minted}{python}
# Nombre de patients par diagnostic
df.groupby("diagnosis")["patient_id"].count()
\end{minted}

% ============================================================
\section{Cr\'eer de nouvelles colonnes}

\begin{minted}{python}
# IMC a partir du poids et de la taille
df["bmi"] = df["weight_kg"] / (df["height_m"] ** 2)

# Groupe d'age
df["age_group"] = pd.cut(df["age"],
                         bins=[0, 18, 40, 60, 100],
                         labels=["<18", "18-39", "40-59", "60+"])

# Indicateur binaire
df["hypertensive"] = (df["systolic_bp"] >= 140).astype(int)
\end{minted}

% ============================================================
\section{Sauvegarder vos r\'esultats}

\begin{minted}{python}
# Sauvegarder en CSV
df.to_csv("cleaned_patients.csv", index=False)

# Sauvegarder en Excel
df.to_excel("cleaned_patients.xlsx", index=False, sheet_name="Patients")
\end{minted}

% ============================================================
\section{Travaux pratiques : analyse de l'esp\'erance de vie (OMS)}

\begin{exercise}
En utilisant le jeu de donn\'ees Gapminder :
\begin{enumerate}
  \item Chargez les donn\'ees et affichez les statistiques de base.
  \item Filtrez les pays africains uniquement. Combien de pays uniques ?
  \item Calculez l'esp\'erance de vie moyenne par d\'ecennie (ann\'ees 1950, 1960, ..., 2000) pour l'Afrique par rapport \`a l'Europe.
  \item Trouvez les 5 pays africains avec l'esp\'erance de vie la plus basse en 2007.
  \item Cr\'eez une colonne \texttt{life\_exp\_category} : <<\,Basse\,>> ($<55$), <<\,Moyenne\,>> ($55$--$70$), <<\,\'Elev\'ee\,>> ($>70$).
  \item Regroupez par continent et \texttt{life\_exp\_category}. Combien de pays dans chaque cat\'egorie ?
  \item Sauvegardez les donn\'ees de l'Afrique uniquement dans un fichier CSV.
\end{enumerate}
\end{exercise}

% ============================================================
\section{R\'esum\'e du chapitre}

\begin{itemize}
  \item Un DataFrame est un tableau. Lignes = observations, colonnes = variables.
  \item \texttt{pd.read\_csv()} charge les donn\'ees ; \texttt{.head()}, \texttt{.describe()}, \texttt{.info()} les explorent.
  \item Filtrez avec des conditions bool\'eennes : \texttt{df[df["col"] > valeur]}.
  \item \texttt{.groupby()} + \texttt{.agg()} r\'epond aux questions <<\,par groupe\,>>.
  \item \texttt{pd.cut()} cr\'ee des variables cat\'egorielles \`a partir de variables continues.
  \item Sauvegardez toujours les donn\'ees nettoy\'ees pour la reproductibilit\'e.
\end{itemize}
